\chapter{Conclusion}

In this thesis we investigated two type systems for Type-GBFS that approximate a type system based on the bench transition system. We implemented both type systems with variants for different type and state selection strategies, and evaluated all configurations on a suite of satisficing benchmarks from IPC 2011 and 2014. Aside from replicating some of the results obtained by \tcite{tomasz2024}, we independently modified three parameters that control the behavior of the algorithms. We observed small improvements in coverage at various type and state selection temperatures as well as exploration rates, but none produced a result that clearly exceeds the findings of the original work. In addition to coverage, we also measured how solution lengths respond to different temperatures and exploration rates. Our results suggest that the two best-performing configurations, \hiDu and \hiDh, suffer from increased solution lengths at the default parameter values, which we attribute to over-aggressive depth biasing: we observed a considerable reduction in solution lengths at increased type selection temperatures and at reduced exploration rates, with no large losses in coverage in either case. Future work is needed to test whether our observations hold for different suites of problems, and whether modifications can be combined to produce stronger improvements. Overall, our results seem to support \tcite{tomasz2024}'s findings that BTS-approximate type systems are a viable method of adding stochastic exploration to GBFS.
